\subsubsection{Lectores enriquecidos y memoria de fibra}
\label{act21:sec:ch-memoria-fibra}

La diferencia entre historia y valor admite un criterio operativo. En un
nivel finito \(H_k\), sea \(q_k:H_k\to Z_k\) un lector y sea
\(M_k:H_k\to\mathcal M_k\) el registro conservado. El lector enriquecido
es \(\widehat q_k=(q_k,M_k)\). La memoria distingue estados de una misma
fibra del lector; su suficiencia se comprueba en el dominio considerado.
Esta formulación incorpora al propio artículo el desarrollo completo de la
información de fibra que se utiliza a continuación.

\begin{proposition}[Recuperabilidad en un nivel finito]
\label{act21:prop:ch-recuperabilidad-fibra}
Existe una inversa de \(\widehat q_k\) sobre su imagen si y sólo si
\(\widehat q_k\) es inyectivo. Para una variable aleatoria \(Z\) con
valores en \(H_k\) y distribución \(\mu_k\), se cumple
\[
 \mathsf H(Z)=\mathsf H(q_k(Z))+\mathsf H(Z\mid q_k(Z)).
\]
Si \(\widehat q_k\) es inyectivo sobre el soporte de \(Z\), entonces
\[
 \mathsf H(Z\mid q_k(Z),M_k(Z))=0.
\]
Aquí \(\mathsf H(Z)=-\sum_z\Pr(Z=z)\log\Pr(Z=z)\), con
\(0\log0=0\), es la entropía finita en nats; la entropía condicional es
el promedio de la entropía de las distribuciones condicionadas.
\end{proposition}

\begin{proof}
Una inversa recupera un único antecedente, lo que exige inyectividad;
recíprocamente, ésta permite asignar a cada par expresado su único
antecedente. Puesto que \(q_k(Z)\) es función determinista de \(Z\),
agrupar la suma que define \(\mathsf H(Z)\) por fibras de \(q_k\) da la
regla de la cadena. Si el par \((q_k(Z),M_k(Z))\) distingue el soporte,
cada distribución condicionada correspondiente está concentrada en un solo
estado y su entropía es cero.
\end{proof}

No expresar una coordenada y borrar el registro que permite recuperarla son
operaciones distintas. Este resultado no identifica entropía con cardinalidad
del continuo ni introduce una interpretación térmica: su dominio, lector y
distribución son los que se acaban de fijar.

\subsubsection{Equivarianza de la medida bajo simetrías del árbol}
\label{act21:sec:ch-naturalidad-medida}

\begin{proposition}[Transporte conjunto del árbol y de su medida]
\label{act21:prop:ch-naturalidad-medida}
Sea \(g=(g_k)_{k\geq0}\) una colección de biyecciones \(g_k:H_k\to H_k\)
que conmuta con los truncamientos, fija la raíz y transporta biyectivamente
las emisiones admisibles, conservando sus multiplicidades. Entonces
\[
 d(g_k\zeta)=d(\zeta),\qquad
 (g_k)_*\mu_k=\mu_k,\qquad (g_\infty)_*\mu=\mu.
\]
Más generalmente, si \(\mu_\zeta\) es la ley de continuaciones de una
historia \(\zeta\in H_k\), se tiene
\[
 (g_\infty)_*\mu_\zeta=\mu_{g_k\zeta}.
\]
\end{proposition}

\begin{proof}
La biyección de emisiones conserva su número con multiplicidades. La
distribución uniforme de semillas también es invariante. Cada trayectoria
y su imagen reciben, por tanto, el mismo producto de pesos locales
\(1/d(\zeta)\). Se obtiene la igualdad de medidas en cada nivel y en
cada cilindro. La coherencia con los truncamientos define \(g_\infty\),
y la unicidad de la medida determinada por cilindros prueba la igualdad
en el límite. El mismo argumento, comenzando en \(\zeta\), prueba la
afirmación sobre continuaciones.
\end{proof}

Se obtiene así la naturalidad del árbol y de la medida: un cambio coherente
de representación transporta los pesos y no obliga a escogerlos de nuevo.
La conservación de emisiones y multiplicidades es una hipótesis esencial de
la proposición, no un dato añadido después de la construcción.
